% Part: applied-modal-logics
% Chapter: epistemic-logic
% Section: properties-accessibility

\documentclass[../../../include/open-logic-section]{subfiles}

\begin{document}

\olfileid{aml}{el}{acc}

\olsection{Relasi Aksesibilitas lan Prinsip Epistemik}

Adhedhasar kang wis kita ngerteni bab korespondensi kerangka ing logika
modal lumrah, kita bisa nliti wujud !!{formula}s khas yen diwenehi
tafsiran epistemik. Kita wis nyebutake yen logika epistemik biasane
ditafsirake ing S5. Saiki ayo dideleng carane maneka prinsip epistemik
diwakili lan carane prinsip iku cocog karo maneka syarat kerangka.

Elinga saka logika modal lumrah yen !!{formula}s modal kang beda-beda
nggambarake sipat relasi aksesibilitas kang beda-beda. Tabel iki
milih sawetara rumus kang cocog karo prinsip epistemik tartamtu.

\begin{table}[t]
    \begin{tabular}{| p{.48\textwidth} || p{.48\textwidth} |}
      \hline
      {\emph{Yen $R$ iku \dots}} & {\emph{mula \dots bener ing~$\mModel{M}$:}} \\
      \hline \hline
        
      & $\Knows (p \lif q) \lif (\Knows p \lif \Knows q)$ 
      \hfill \newline{(Katutupan)} \\
      \hline
      \emph{refleksif}: $\forall w Rww$  
      & $\Knows p \lif p$ \hfill \newline{(Veridikalitas)} \\
      \hline
      \emph{transitif}: \newline
      $\forall u \forall v \forall w ((Ruv \land Rvw) \lif Ruw)$ & 
      $\Knows p \lif \Knows \Knows p$ \hfill 
           \newline{(Introspeksi Positif)} \\
      \hline 
      \emph{Euklidean}: \newline
      $\forall w \forall u \forall v ((Rwu \land Rwv) \lif Ruv)$ &  
      $\lnot \Knows p \lif \Knows \neg \Knows p$ \hfill 
        \newline{(Introspeksi Negatif)} \\
      \hline
    \end{tabular}
    \caption{Papat prinsip epistemik.}
    \ollabel{tab:four}
  \end{table} 

Veridikalitas, kang cocog karo aksioma~$T$, kerep dianggep prinsip
kang paling ora dadi pasulayan, amarga nyatakake yen !!a{formula}
dingerteni, rumus iku mesthi bener. Katutupan, uga introspeksi positif
lan negatif, luwih akeh dadi pasulayan.

Katutupan, kang cocog karo aksioma~$K$, nglambangake gagasan yen
kawruh sawijining agen katutup miturut implikasi. Ing sawetara
kahanan, iki katon lumrah. Contone, aku bisa ngerti manawa yen aku ana
ing Victoria, mula aku ana ing Pulo Vancouver. Yen ora ana skenario skeptis
kang aneh, aku pancen ngerti yen aku ana ing Victoria, mula iki
uga nuduhake yen aku ngerti yen aku ana ing Pulo Vancouver. Dadi,
ing kasus iki, katutupan logis kawruhku katon cukup lumrah.
Nanging, kita ora tansah mikirake sakabehe konsekuensi saka kang
kita ngerteni, mula ing kasus liyane prinsip iki bisa marakake
asil kang ora patiya lumrah.

Introspeksi positif, kang kadhang diarani prinsip KK, bisa diandharake
kanthi pratelan: yen aku ngerti sawijining bab, aku ngerti yen aku
ngerti bab iku. Iki pasangan epistemik aksioma~4. Kajaba iku,
introspeksi negatif bisa diandharake kanthi pratelan: yen aku
\emph{ora} ngerti sawijining bab, aku ngerti yen aku ora ngerti
bab iku. Iki pasangan epistemik aksioma~5. Loro prinsip iki katon
nduweni tuladha kang mbantah kanthi cukup lumrah: aku bisa ora yakin
apa aku ngerti sawijining bab, senadyan nyatane aku ngerti bab iku.


\end{document}
